\originalchapter{8.322作业十二与逐题未取边界}
\originalsection{Q1：KG台阶的全部能区}
\begin{exercise}
一维$V=0$（$x<0$）、$V=V_0$（$x>0$），左入射Klein--Gordon平面波。求反射与透射系数，讨论不同能区。
\end{exercise}
\begin{solution}
设$E>mc^2$、$k=\sqrt{E^2-m^2c^4}/(\hbar c)>0$。固定频率场为左$e^{ikx}+re^{-ikx}$、右$te^{iqx}$，$q^2=[(E-V_0)^2-m^2c^4]/(\hbar^2c^2)$。因静态KG方程无delta项，$\phi,\phi'$连续，$r=(k-q)/(k+q)$、$t=2k/(k+q)$。流$\propto\operatorname{Im}\phi^*\phi'$给
\[
 R=|r|^2,\qquad T=\frac qk|t|^2=\frac{4kq}{(k+q)^2},\quad R+T=1
\]
用于实$q$，这里$T$是电荷流比。

若$E-V_0>mc^2$，选$q>0$，群速正、$R,T$均在0--1。若$|E-V_0|<mc^2$，选$q=i\kappa$衰减，右无通往无穷的流，$R=1,T=0$。若$E-V_0<-mc^2$，为了右向群速$\partial E/\partial(\hbar q)=\hbar c^2q/(E-V_0)>0$，须选$q<0$，故$T<0,R>1$，这是Klein区。不能将负透射当单粒子负概率，需反粒子与外场供能的场论解释。若错误选$q>0$，所得正透射实际上右侧群速度指向界面，属于额外右入射边界。在孤立的$k+q=0$处这些驻定匹配系数有极点，不能声称有限单粒子散射概率存在。台阶无限延伸且强场的形式单粒子模型只给流关系，不解决完整粒子产生过程。
\end{solution}

\originalsection{Q2：按指定矩阵的Weyl表示}
\begin{exercise}
给$U=\frac1{\sqrt2}\left(\begin{smallmatrix}1&1\\-1&1\end{smallmatrix}\right)$（每格为二维单位矩阵），$\gamma_W^\mu=U\gamma_D^\mu U^{-1}$。求新矩阵，写$\psi=(\phi,\chi)^T$的Dirac方程及$m\to0$极限。
\end{exercise}
\begin{solution}
逐块乘法给
\[
 \gamma_W^0=\begin{pmatrix}0&-I\\-I&0\end{pmatrix},\quad
 \gamma_W^i=\begin{pmatrix}0&\sigma_i\\-\sigma_i&0\end{pmatrix},\quad
 \gamma_W^5=\begin{pmatrix}I&0\\0&-I\end{pmatrix}.
\]
注意指定$U$使$\gamma^0$带负号；不能背诵另一常见Weyl约定的正号。取$\hbar=c=1$，方程等价于
$i(\partial_t+\sigma\cdot\nabla)\phi=-m\chi$、
$i(\partial_t-\sigma\cdot\nabla)\chi=-m\phi$。
$m=0$时两条Weyl方程独立，$\phi$是$\gamma^5=+1$、$\chi$是$-1$的手征分量；正能平面波分别有正、负helicity。质量项耦合两手征，不能对非零质量任意忽略其中一个而仍满足Dirac方程。
\end{solution}

\originalsection{Q3：16个Clifford积的独立性}
\begin{exercise}
四个$\gamma^\mu$每个取或不取、按固定顺序相乘的16个积$\Gamma^a$，证明：(a)平方为$\pm I$；(b)除单位外迹0；(c)$\Gamma^a\Gamma^b$等于另一积乘符号，仅$a=b$时为单位；据此证明线性独立。
\end{exercise}
\begin{solution}
积对应指标子集$A\subset\{0,1,2,3\}$。不同gamma反对易，$\gamma_0^2=I,\gamma_i^2=-I$，故把同指标移相邻消去，$\Gamma_A^2=(-1)^{|A|(|A|-1)/2}\prod_{\mu\in A}\eta^{\mu\mu}\,I$，验证(a)。任何非空$A$存在一个gamma与积反对易：$|A|$偶时取$A$内指标，奇时取$A$外指标（奇数最大3，外指标存在）。共轭保持迹却给负积，故迹0。乘两积时公共指标平方消去，剩对称差$A\triangle B$及符号，只有$A=B$为空，证明(c)。若$\sum_Ac_A\Gamma_A=0$，左乘$\Gamma_B$取迹，交叉项全0，本项为$\pm4c_B$，所以全部$c_B=0$。16维恰是$M_4(\mathbb C)$维数；平方符号已显式写出，不只称“显然独立”。
\end{solution}

\originalsection{Q4：负能静止态的boost和自旋读数}
\begin{exercise}
初负能静止解$\psi=(0,0,1,0)^Te^{imc^2t/\hbar}$，$S_z=+\hbar/2$。变换到沿原$x$方向以$v$运动的新系，求$\psi'(x')$、能量/动量和新$z'$测自旋$-\hbar/2$的概率。
\end{exercise}
\begin{solution}
设$\tanh\eta=v/c$。被动Lorentz变换给$E'=-\gamma mc^2$、$p'_x=+\gamma mv$。坐标反变换$t=\gamma(t'+vx'/c^2)$给原相位与之匹配。旋or变换$S=e^{-\eta\alpha_x/2}$，故
\[
 \psi'(x')=\begin{pmatrix}0\\-\sinh(\eta/2)\\\cosh(\eta/2)\\0\end{pmatrix}
 e^{i\gamma mc^2t'/\hbar+i\gamma mvx'/\hbar}.
\]
在定常平面波的四分量$\Sigma_z/2$读数中，负自旋仅第二分量，归一分母$\cosh^2(\eta/2)+\sinh^2(\eta/2)=\gamma$，所以$P_-=(\gamma-1)/(2\gamma)$、$P_+=(\gamma+1)/(2\gamma)$。$v\to0$回0，极高速趋1/2。这个问题明确测Dirac矩阵自旋分量；相对论静止自旋/Pauli--Lubanski投影的实验定义未必等同$\Sigma_z/2$，不能由此结论说粒子的固有自旋量子数改变。
\end{solution}

\originalsection{12套作业的逐项状态与缺失保留}
\begin{longtable}{p{.10\textwidth}p{.35\textwidth}p{.45\textwidth}}
\toprule 作业 & 本次完整题面与自主解答 & 未取得/依赖缺失，逐项保留\\\midrule\endhead
1 & Q5 & Q1--4：Sakurai5.22、5.28、5.29、5.30。\\
2 & 无 & Q1--5：5.33、5.35、5.37、5.38、5.39。\\
3 & Q1、Q3--5 & Q2：5.40。其Q3完整给具体任务，已从头补基础速率，不借未附答案。\\
4 & Q2--3 & Q1官方因版权撤除；Q4引用Sakurai S.14的$\vct V_m(R)$定义未附，未冒称完整原题已解。上章仅给自定义Berry曲率的独立无散度说明。\\
5 & 无 & Q1--5：4.3、4.5、4.6、4.8、4.9。\\
6 & Q3--4 & Q1--2：4.11、4.12。\\
7 & Q5--6 & Q1--4：6.1、6.2、6.4、6.5。\\
8 & Q2--6 & Q1：6.6。\\
9 & Q1、Q3 & Q2：7.1。\\
10 & Q2 & Q1：7.2；Q3--5：7.3、7.4、7.7。\\
11 & Q1、Q4--5 & Q2--3：7.8、7.9。\\
12 & Q1--4 & 四题实际题面均可取得。\\\bottomrule\end{longtable}
合计26道实际可用完整大题自写解答，29道纯书号、1道撤除、1道定义依赖未取，原编号共57项。与8.321新增20道纯书号和旧本科10道缺项分开，未把书号当作新完成题。没有取得8.322官方试卷或解答，因此所有新增解答的身份是自主解答。编译和图示只验证表达/数值复现，不替代严谨推导，新增公式检查见研究生QA记录。
